\chapter{Functional Scope}
\label{chap:functional-scope}

This chapter outlines the considerations and requirements for the library and the supporting tools. Describes the intended audience and details the key decisions made during development. These choices are further described in subsequent sections on Functional and Non-functional requirements.

\section{Expected User Groups}
Due to the experimental nature of this work, the key design goal was to make the library both easy to use and easy to understand. This is necessary as we anticipate that users may need to adapt the library for integration into their specific applications.

\paragraph{Embedded developers}
The library was designed with the limitations of embedded hardware in mind. Unlike many alternatives, it supports every step of the process on-device, while others implement predictions on-device only. Thanks to the designed model exchange format, users can choose to train models on faster hardware or use the \GreenTsetlinName library to later move them to the edge device. Additionally, the efficiency gains promised by TMs~\cite{granmo2021tsetlinmachinegame} are highly beneficial for these resource limited devices, where standard machine learning methods are often too expensive or complicated to run.

\paragraph{Researchers}
As TMs have only recently started to gain recognition, we expect researchers familiar with classic machine learning models to benefit from integrating them into their workflows. By offering a familiar interface, our library enables rapid prototyping. This distinguishes our solution from others that require deeper domain knowledge, making it easier to benchmark the models in various use cases.

\section{Functional Requirements}

\paragraph{On-device training and inference}
Our solution was designed to support both TM training and inference directly on target hardware. Unlike other embedded machine learning frameworks that perform only inference on edge devices, this library enables the model to learn directly from newly gathered data. This is necessary for applications that adapt to local sensor data without a constant connection to the central server.
It should however facilitate an inference only mode with a lighter model serialization format specifically for more resource constrained systems.

\paragraph{Interoperability and serialization}
The library must implement functionality for saving and loading model states. This requires a standardized serialization format to ensure compatibility with the \GreenTsetlinName library. This compatibility is essential for the users to leverage high-performance workstations to pre-train models, export them, and then deploy them directly to the embedded devices. Furthermore, models trained or tuned on the edge must also be exportable for backup or analysis on faster hardware.

\paragraph{Configurable hyperparameters}
Exposing TM model hyperparameters is necessary for fitting the model to its application with varied limitations dictated by hardware or the input data itself. Researchers and developers must be able to easily fine-tune these parameters when designing their solution for each specific problem. This allows optimization of performance, memory footprint, or accuracy without modifying library code. 

\paragraph{Correctness with the theoretical work}
The implementation must behave as described in theoretical research~\cite{granmo2021tsetlinmachinegame}. It is required to accurately model the Tsetlin Automata logic and implement the specific learning rules (Type I and Type II feedback) defined in the original papers. This ensures that the software is a close representation of the method proposed in the theory and avoids unintended assumptions.

\section{Nonfunctional Requirements}

\paragraph{Portability}
The library must be designed for compatibility, avoiding dependencies on any specific hardware or operating system. The functionality of the external library should be defined during compilation if there is a risk that it may not work on the desired target. This strategy ensures that the code can be compiled with minimal, if any, modifications for a wide range of hardware architectures. The resulting portability, coupled with the adjustable model parameters, allows end users to create a solution that is both widely adoptable and future-proof.

\paragraph{Performance and memory usage}
Optimizations are significantly beneficial, as the library's performance impacts both the training step and, more importantly, the real-time predictions on target devices. Given that TMs heavily rely on binary logic and random number generators, optimizing these areas allows for better utilization of limited hardware resources. This optimization extends to memory management, since the memory is needed to store the model on the device, as well as the memory used during runtime, must also be considered.

\paragraph{Simple API}
The Applications Programming Interface (API) must be intuitive and well-documented. Its design should follow standard interfaces developed by major machine learning libraries, allowing users unfamiliar with TM theory to quickly prototype concepts. This simple interface allows a developer to easily integrate the library into an existing project. Moreover, supporting quick prototyping in environments like Python Jupyter Notebooks alongside established models from libraries such as Scikit-learn speeds up development before deployment to embedded devices.

\paragraph{Testing and examples}
As with any software project, testing enhances the development process and maintenance. Code examples provide a starting point for users unfamiliar with the project. Furthermore, simple installation and easy to run examples encourage potential users to quickly try the library. Finally, these tests help verify that the library runs correctly on new devices with different hardware limitations.